% R014_BOUNDED_PREDICTION_REVIEW
\subsubsection{Predictions as Bounded Research Targets}
\label{sec:predictions-bounded-r014}

The prediction chapters in Core 1.3.3 are not promoted as a completed
cross-domain empirical theory. They are a review surface for asking whether a
candidate statement has a declared level, a formal vocabulary, an evidence row,
a comparator, and a falsifier. A row may therefore be read as a target for
future test design only when it names the assumptions under which the target is
meant to operate.

The release uses the following conservative rule. A candidate prediction at
level \(K_x\) is public-supporting only if it preserves the tuple components
\[
  (\Omega,\partial\Omega,A,\Theta,P,J,C,k,M)
\]
or explicitly states which component is derived in the local model. It must also
state what would reopen the row: a missing source snapshot, a failed finite
witness, a successful negative control, a stronger comparator, or a domain
observation that requires variables outside the declared tuple.

Consequently, the K-level prediction sections below should be read as bounded
review checklists. They do not claim unrestricted prediction across physics,
chemistry, biology, cognition, society, civilization, or meta-theory. They say
what kind of evidence would be needed before a level-specific statement could be
promoted beyond illustrative or programmatic status.

\paragraph{Review checklist.}
The reviewer should ask six questions before treating any row as public
support. First, is the target level \(K_x\) named with enough precision that an
adjacent level would make a different claim? Second, is the tuple component
being tested visible in the prose or formula? Third, is the source family known:
proof sheet, finite witness, replay table, comparator register, figure, table,
or domain note? Fourth, does the row identify what a negative control would
look like? Fifth, is the comparator boundary strong enough to prevent the row
from merely renaming a standard account? Sixth, does the text say what finding
would force demotion?

These questions matter because a prediction surface can fail in several
different ways. It can be formally named but evidentially empty. It can be
empirically suggestive but too broad for the tuple. It can be clear inside one
domain and misleading when copied into another. It can be useful as a design
heuristic but not yet usable as a scientific result. Core 1.3.3 keeps those
states separate so that the reader can criticize the release without guessing
whether an example, theorem route, replay row, or future research target is
being used.

\paragraph{Promotion boundary.}
A prediction row moves upward only when the support moves upward. A bounded
example remains an example. A formula without source data remains a formal
anchor. A replay row without comparator and negative control remains a replay
exercise. A proof-sheet reference remains a proof/evidence anchor inside its
assumptions, not an empirical closure claim about a whole domain. This
separation is deliberately conservative: it makes the current release less
spectacular, but more inspectable.

\paragraph{Prediction review taxonomy.}
The release uses five kinds of prediction-adjacent material. The first kind is
a theorem-route constraint. It says that if a reader accepts the declared
definitions and assumptions, a stated obstruction, boundary, identity, or
continuumness result follows inside the formal route. This kind of row is
reviewed by checking definitions, proof-sheet dependencies, finite witnesses,
and assumption drift. It is not reviewed by asking whether the same sentence is
already source-bound by completed domain evidence.

The second kind is a finite or computational witness. It says that a small
model, countermodel, keep/drop pair, or semantic check illustrates that the
formal condition is non-vacuous or that a rejected inference fails. This kind of
row is useful because it prevents purely verbal closure, but it is still local:
it supports the existence or non-existence of a pattern under the declared
encoding rather than a broad empirical law.

The third kind is a replay row. A replay row connects a source snapshot, a
formula or reconstruction rule, an observed value or expected finite result, a
comparator, a residual or uncertainty statement, a negative control, and a
falsifier. This is the closest current surface to empirical testing, but it is
bounded by its source and method. A replay row does not validate a whole domain
merely because the command or table exists.

The fourth kind is a comparator-boundary row. It asks whether a standard
framework, predecessor theory, domain method, or simpler model already explains
the same target with less burden. This row is especially important for a theory
with a broad title. It keeps novelty language modest and forces OC to say what
residual work the tuple, K-levels, proof/evidence discipline, or reopening rule
actually add.

The fifth kind is an illustrative or planned-work row. It may be useful for
teaching, domain orientation, or future experiment design, but it should not be
quoted as promoted support. Its value is that it shows where a better source
snapshot, formula, comparator, or falsifier would be needed. In Core 1.3.3, a
large part of the high-level social, civilizational, and meta-theoretical
material is intentionally handled this way unless the specific row carries its
own support package.

\paragraph{How to criticize these rows.}
A reviewer can criticize a prediction row without accepting OC's larger
ambition. The strongest criticism is not "this sounds broad" but "this exact row
has no source, no formal anchor, no comparator, no negative control, or no
reopening condition." The release is designed so that such a criticism can be
localized. If the problem is formal, the theorem route is reopened. If the
problem is empirical, the replay row is demoted. If the problem is novelty, the
comparator boundary is narrowed. If the problem is didactic, the example is
rewritten without changing the scientific claim.

This localization is part of the scientific method of the release. It prevents
one weak example from silently invalidating the entire tuple, and it prevents
one formal proof from silently licensing an entire empirical domain. The
manuscript is therefore deliberately redundant in its claim surfaces: prose,
formula, figure, table, proof sheet, replay row, comparator row, and falsifier
are different review instruments.
